\newpage
\topic{23 continued}{Guided reading and video lectures}
\textbf{1. From momentum to a pressure equation.}
Chris Greenshields and Henry Weller, \emph{Notes on Computational Fluid Dynamics: General Principles},
\href{https://doc.cfd.direct/notes/cfd-general-principles/pressure-velocity-coupling}{section 5.10, Pressure--velocity coupling}.
Read alongside Sections 4--7 of this lecture. Follow the momentum diagonal,
the face-flux correction and substitution into continuity. Their software
notation differs from the staggered scalar equations used here.

\textbf{2. SIMPLE and under-relaxation.}
The same authors' \href{https://doc.cfd.direct/notes/cfd-general-principles/steady-state-solution}{section 5.12, Steady-state solution},
is a companion to Section 8. Identify the momentum predictor, conservative
flux correction and next nonlinear iteration. Compare it with Aidan Wimshurst's
Fluid Mechanics 101 video,
\href{https://www.youtube.com/watch?v=OOILoJ1zuiw}{\emph{[CFD] The SIMPLE Algorithm (to solve incompressible Navier--Stokes)}}.
Use the diagram in this lecture to track where each unknown is stored while
following the pressure-correction sequence in the video.

\textbf{3. PISO within a physical time step.}
Read \href{https://doc.cfd.direct/notes/cfd-general-principles/transient-solution}{section 5.19, Transient solution},
alongside Sections 9--11. Then watch Fluid Mechanics 101,
\href{https://www.youtube.com/watch?v=ahdW5TKacok}{\emph{[CFD] The PISO Algorithm}}.
Focus on what is updated before the next pressure solve, and what remains
fixed during the time step. Non-orthogonal corrections discussed in the video
are a separate mesh issue; the present uniform Cartesian grid does not require them.

\textbf{4. PIMPLE and nonlinear outer iterations.}
Read \href{https://doc.cfd.direct/notes/cfd-general-principles/the-pimple-algorithm}{section 5.21, The PIMPLE algorithm},
alongside Sections 12--14. Compare the outer coefficient updates with the
inner corrections. Distinguish time-accurate integration from a
pseudo-transient route to a steady state; more correctors do not replace a
time-step refinement study.

\textbf{5. Connect the loops to practical controls.}
The official \href{https://doc.cfd.direct/openfoam/user-guide-v13/fvsolution}{OpenFOAM v13 User Guide, Solution and algorithm control},
explains pressure references, \texttt{nCorrectors} and
\texttt{nOuterCorrectors}. Map these names to the loops before reading
software configuration examples. This Python lesson uses its own MAC-grid
implementation, so software-specific interpolation details and defaults
should not be transferred without checking.

The videos are linked as companion explanations. The
\href{https://www.fluidmechanics101.com/pages/lectures.html}{author's lecture catalogue}
lists the SIMPLE and PISO lectures and distinguishes the YouTube videos
from separately sold slide collections. No paid slides are required for this lesson.
